\section{Introduction}
\label{sec:intro}
% [2 pts] Briefly describe and summarize the goal of this assignment.

Ordinary diffusion sampling starts from pure noise and runs the reverse process step by step to produce an image.
Score distillation instead treats the image (here, a Stable Diffusion latent) as a set of free parameters and uses a frozen, pretrained diffusion model only as a \emph{critic}: at every iteration the model looks at a noisy version of the current latent and tells us in which direction it should move to look more like the text prompt.
This turns a diffusion model into a loss function, which is what makes text-to-3D methods such as DreamFusion~\cite{poole2023dreamfusion} possible (\cref{fig:teaser}).

In this assignment I implemented three score distillation methods on top of Stable Diffusion 2.1-base for text-to-image generation:
Score Distillation Sampling (SDS)~\cite{poole2023dreamfusion},
Variational Score Distillation (VSD)~\cite{wang2023prolificdreamer}, and
Score Distillation via Inversion (SDI)~\cite{lukoianov2024sdi}.
Each method optimizes a $1\times4\times64\times64$ latent for the 10 evaluation prompts, and the results are evaluated with the CLIP ViT-B/32 image--text similarity~\cite{radford2021clip}.
I also ran the required ablations on the SDS guidance scale and on the SDI inversion steps and update interval, and an extra experiment with the original diffusion loss.

\Cref{sec:method} describes the methods and my implementation, \cref{sec:exp} reports the results and ablations, and \cref{sec:discussion} discusses why the plain diffusion loss fails and why VSD and SDI improve over SDS.
